% ======================================================================
\section{Experimental Protocol}
\label{sec:setup}

\subsection{Corpus}
\tabv{tab_datasets_v2}
Eight national systems---Germany, France, Great Britain, Spain, Italy, the Netherlands,
Poland and Belgium---are taken from the Open Power System Data time-series package
\cite{opsd2020}, which carries ENTSO-E actual load and, for every system, the transmission
operator's own published day-ahead load forecast; they are joined to ERA5 reanalysis weather
\cite{hersbach2020era5}, obtained through the Open-Meteo historical-weather interface
\cite{openmeteo} at one interior point per system near its geographic centre (2\,m
temperature, relative humidity, shortwave radiation and 10\,m wind speed), and to solar and
wind generation where present. A single point is a simplification: national load responds to
population-weighted weather, and the two covariate pathways compared in this paper need not
profit equally from a richer representation, so Section~\ref{sec:sensitivity} repeats the
comparison of the two leading models with the population-weighted mean of six metropolitan
points on three of the larger systems, France, Spain and Italy, where one point is least
representative of the national weather. Each system carries local civil timestamps and its national
public holidays from the \texttt{holidays} package (for Great Britain the England calendar,
which covers most of the load). The eight systems were fixed before any experiment: the three
largest western European systems and five further systems chosen for a spread of size,
climate and holiday calendar; other systems in the package carry the same columns and can be
added under the protocol unchanged. A generation series is dropped by two rules fixed in
advance: any single internal gap longer than thirty days, or less than half of the load span
present (Table~\ref{tab:datasets-v2}). Remaining gaps are interpolated linearly up to six
hours and otherwise filled forward and backward; the counts are recorded with the data.
Three features of the published series matter for reading the results and are documented with
the data. The Dutch actual-load series steps up in January 2016 (mean load
\nv{data.OPSD-NL.step_pct} higher in 2016 than in 2015, a change of reporting perimeter inside
the training split), which any model with a linear trend extrapolates. The French series carries
one physically impossible hour in the test split (\nv{data.OPSD-FR.bad_hour.value_gw}\,GW at
\nv{data.OPSD-FR.bad_hour.stamp} between neighbours of
\nv{data.OPSD-FR.bad_hour.lo_gw}--\nv{data.OPSD-FR.bad_hour.hi_gw}\,GW); the windows whose
targets cover it are excluded from every French point score, the hour from the probabilistic
scores and the day from the day-ahead protocol, and the record keeps the unmasked values. The
British series has \nv{data.OPSD-GB.n_spikes} isolated single-hour spikes (an hour that departs
from the mean of its neighbours by more than \nv{data.spike.rel_pct} of mean load while the
neighbours agree within \nv{data.spike.neigh_pct}; the other systems have at most a handful),
which are left in place and documented with the data.
Splits are chronological, 70/10/20, so that the test span (August 2019 to September 2020)
contains one year-end holiday period and the spring-2020 lockdowns; the validation span
contains neither. For comparability with the literature we also report the 321-client
Electricity benchmark \cite{lai2018lstnet}, rebuilt from the UCI record with true timestamps
\cite{uciecl}, and the four ETT transformer-station series \cite{zhou2021informer}, on which
the calendar prior is expected to fail.

\subsection{Forecasters}
\label{sec:baselines}
\nv{n.models.learned} learned forecasters are compared: \nv{proto.n_neural} neural models, the
seasonal-na\"ive reference among them, and two classical references (gradient-boosted trees,
and Hong's regression in two fitting schemes that count as one model, plus the same regression
with an autoregressive model of its residual, which counts as another). The neural models pass
through one training loop on the same windows, normalisation, optimiser, budget and inputs
(the seasonal-na\"ive reference has nothing to train; only its period is chosen on
validation): \nv{proto.n_v1_models} compact re-implementations following the published
architectures---Seasonal-Na\"ive; LSTM \cite{kong2019lstm} and TCN \cite{bai2018tcn}; N-BEATS
\cite{oreshkin2020nbeats}; NLinear and DLinear \cite{zeng2023dlinear}, RLinear
\cite{li2023rlinear} and FITS \cite{xu2024fits}; TimeMixer \cite{wang2024timemixer}; Informer
\cite{zhou2021informer}, Autoformer \cite{wu2021autoformer}, PatchTST \cite{nie2023patchtst}
and iTransformer \cite{liu2024itransformer}; the calendar-conditioned control PatchTST+Cal,
which is handed PLANET's design matrix as tokens without the solver; PLANET; and its
native-covariate variant PLANET-X (Section~\ref{sec:planetx})---and
\nv{proto.n_nf_models} taken from the NeuralForecast library \cite{neuralforecast}: the
Temporal Fusion Transformer \cite{lim2021tft}, TiDE \cite{das2023tide}, N-HiTS
\cite{challu2023nhits}, N-BEATSx \cite{olivares2023nbeatsx},
TimeXer \cite{wang2024timexer}, TSMixerx \cite{chen2023tsmixer} and XLinear
\cite{xlinear2026}. Two classical references complete the set: gradient-boosted trees on lag,
calendar and weather features \cite{friedman2001gbm}, and Hong's vanilla regression with the
recency extension \cite{hong2016gefcom,wang2016recency}, reported fitted once, re-fitted
monthly (the fitting window, expanding or the last 36 or 24 months, and the linear trend
chosen on validation), and re-fitted monthly with an autoregressive model of the residual: the
regression with AR errors of the forecasting literature \cite{hyndman2021fpp3}, whose AR
coefficients are estimated once on the training residuals and whose order is chosen on
validation among 0, 24, 48, 168 and 240 hours (240 hours of residuals plus the 72 hours of
recency lags they need is the most the common lookback of 336 hours admits; the bound is
selected on most systems, and the validation loss changes little between 168 and 240 hours),
so that the classical reference also has access to the recent load history that every neural
model reads. Two zero-shot foundation models, Chronos-2
\cite{chronos2} and TimesFM-2.0 \cite{das2024timesfm}, are run without training as reference
points, with the caveat that their pre-training corpora may overlap the public series used
here.

Every model sees per-window instance-normalised inputs: reversible instance normalisation in
the re-implementations, PLANET's own normalisation, and, for the library models, which the
harness calls directly rather than through the library's training loop, the same per-window
standardisation of the lookback applied in the wrapper and inverted on the forecast. Two
re-implementations depart from this: Informer and Autoformer run without instance
normalisation, on the standardised block as their original architectures do, because the
re-run with it added that the protocol called for did not fit the compute budget of the study;
their cells should be read with that caveat, and the literature on the technique
\cite{kim2022revin} suggests that it costs the two models accuracy rather than lends them any. The
library models run at their published defaults with one exception: TiDE's layer
normalisation is disabled, because the library applies it to the one-dimensional output of
the temporal decoder, which collapses point forecasts to a constant. Under the regimes of
Section~\ref{sec:regimes}, TFT, TiDE, N-HiTS, N-BEATSx, TSMixerx and XLinear consume the
covariate block natively; TimeXer natively reads only its
past part and receives the future part through the shared head; every other learned model
receives the whole block through the head. One control variant isolates the pathway itself:
TiDE run with its native covariate pathway switched off and the block routed through the shared
head at TiDE's own selected settings (TiDE via head, Section~\ref{sec:mech}). Two architectures
carry calendar information of their own in every regime, and the reader should know it: PLANET
and PatchTST+Cal through the
design matrix, and Informer and Autoformer through the timestamp embedding that is part of
their architecture (hour, weekday, day, month and day of year, in UTC). The two classical
references consume the subset of the block that their specification admits: Hong's regression
the temperature only, the trees the lag-one past values and the values at the target step.

Every model is run under the headline regime R2. Ten \emph{headline models}---N-BEATS,
PatchTST, PatchTST+Cal, PLANET, PLANET-X, TFT, TiDE, N-HiTS, TimeXer and DLinear, chosen to
span the architectural families---together with the trees are also run under R0 and R1 (Hong's
regression cannot run without a calendar), and the zero-shot models under R0 and R2.

\subsection{Windows, tuning and budget}
Windows are indexed by their first lookback position; the lookback may reach into the previous
split while targets stay inside it; training windows are subsampled with a stride of five
hours on the grids (seven on Electricity, one on the hourly ETT series), coprime with the
daily period so that every start hour is represented, and validation and test windows are
taken at unit stride. Every neural model is trained with Adam under a cosine schedule to early
stopping on validation loss (patience five, cap thirty epochs; the runs stopped after
\nv{proto.epochs.mean} epochs on average and \nv{proto.epochs.frac_capped} of them reached the
cap; Section~\ref{sec:discussion} reports a three-fold budget check), with mean squared error
on the load channel as the training loss, at batch size 64 (16 on Electricity and for
the four most expensive models at the longest lookbacks and horizon, halved after a memory
failure).

Tuning follows one budget for every model, the expensive ones included. On the two reference
systems (DE, FR) every model with a width setting is trained at three widths---half,
published and double its published size; the one-layer linear models and the
seasonal-na\"ive reference have none---and five learning rates spanning two decades,
$\{10^{-4},5\!\times\!10^{-4},10^{-3},3\!\times\!10^{-3},10^{-2}\}$; the width with the
lowest mean validation loss is transferred to the other six systems, where the learning rate
is re-selected from the same five values. PLANET and PLANET-X carry one further tuned
setting, the residual backbone, searched on the reference systems jointly with the width
(twice the configurations of any other model) and transferred; the ablation of
Section~\ref{sec:ablation} reports both backbones. The lookback is then selected on validation
from $\{168,336,512,720\}$ for every model on every system and horizon; the Electricity and ETT
benchmarks and the day-ahead protocol use the protocol lookback of 336. \appref{app:config}
lists every selection and how often it fell on an edge of a grid. The regimes R0 and R1 reuse
the rate, width and lookback selected under R2; re-selecting the learning rate under R0 and R1
on Germany changes the choice for \nv{lrsel.R0.n_changed} and \nv{lrsel.R1.n_changed} of the
\nv{lrsel.R0.n_models} headline models and changes their validation loss by
\nv{lrsel.R0.mean_val_dmse_pct} and \nv{lrsel.R1.mean_val_dmse_pct} on average, so the regime
comparison is read with that margin in mind, and the cells concerned are re-run with their own
rate as a sensitivity (Section~\ref{sec:sensitivity}). On Electricity and ETT the width is the one chosen
on the reference systems and the learning rate is re-selected on each dataset's own validation
split. The four most expensive models---TFT, TimeXer,
Informer and Autoformer---cost between \nv{cost.expensive.ratio_min} and
\nv{cost.expensive.ratio_max} times a PatchTST run (Table~\ref{tab:cost-v2}); they receive
the same budget as every other model. TimesNet \cite{wu2023timesnet} is not in the set: under
the same protocol each of its runs holds the card at full load for between one and three hours
(two-dimensional convolutions over the period-folded input at hourly resolution), and the full
tuning, lookback and headline schedule it would need lay beyond the compute budget of the
study. Every selection is made on validation loss, never on
test, and every reported cell is the mean of three training seeds, and the headline, ablation and
benchmark tables print the seed standard deviation beside it. One rule, applied wherever seeds are
averaged, excludes a seed whose training diverged at the selected learning rate: a validation
loss that is non-finite or more than ten times the median of the other seeds of the cell. The
cell then reports the remaining seeds and their number; this happened \nv{proto.n_diverged}
times (\nv{proto.diverged_runs}). A run that fails twice on the 8\,GB card even at the smallest
batch size is reported as an empty cell rather than re-run on other hardware.


\subsection{Day-ahead operator protocol}
\label{sec:dayahead-proto}
Operators issue their day-ahead forecast before the noon market gate. To compare like with
like, learned models are evaluated from windows whose last observed hour is 11:00 UTC on day
$D-1$ and scored over 00:00--23:00 UTC of day $D$ (lead times 13--36\,h), and the operator's
published forecast is scored on exactly the same target hours of the same days. Under regime
R2 both sides then hold a day-ahead weather forecast. The learned models are 48-hour point models trained under the headline protocol (protocol
lookback 336, tuned width and learning rate),
trained once at the end of the training split and not refitted during the test period, which
favours the operators; two complements measure that handicap: a seed ensemble (the mean of the
three seed forecasts, the cheapest operational improvement) and, for PLANET and TiDE, a monthly
refit from scratch on every row before the start of each test month with early stopping on the
two preceding months (one seed). Because publication deadlines differ, the whole protocol is
repeated with the cut-off at 08:00 UTC (lead times 16--39\,h). Every learned forecast is tested
against the operator by a Diebold--Mariano test on the daily mean absolute errors (HAC variance,
lag $n^{1/3}$), and the scores are also computed inside and outside the first-wave lockdown window (10 March
to 31 May 2020) and on holiday days (Section~\ref{sec:dayahead}); the released summaries carry
them by month as well. For two systems the
published operator series sits well below the published actual load on the scored days, by
\nv{da.OPSD-NL.TSO.bias_pct} for the Netherlands and \nv{da.OPSD-GB.TSO.bias_pct} for Great
Britain, an offset far beyond any forecast error that reflects a difference of perimeter
between the two ENTSO-E series and that varies over time, so it cannot be removed by a
correction estimated on earlier data. The operator is therefore reported as published, with
its bias beside it, and the comparison with the operators is read on the six systems whose
published series is unbiased on the scored days (absolute bias at most \nv{da.grid.bias_tol_pct}
of mean load, a threshold fixed after the biases were seen; it separates two series at
\nv{da.OPSD-NL.TSO.bias_pct} and \nv{da.OPSD-GB.TSO.bias_pct} from six within one per cent).

\subsection{Probabilistic protocol}
\label{sec:prob}
Nine headline models with a quantile head (PLANET, PLANET-X, N-BEATS, PatchTST,
PatchTST+Cal, DLinear, TiDE, N-HiTS and TFT) emit seven quantile levels (5, 10, 25, 50, 75, 90
and 95\,\%) trained with the pinball loss at the tuned width, learning rate, backbone and
lookback of the point model, over three seeds, for the day-ahead horizon. Each symmetric pair
of levels is calibrated by conformalised quantile regression on the validation split, per
horizon step \cite{romano2019cqr}, and the calibrated quantiles are rearranged to be monotone.
The 90\,\% interval is then adjusted online across the test split by adaptive conformal
inference \cite{gibbs2021aci} in its delayed-feedback form: the working level of horizon step
$k$ at a given origin is updated only with the outcome of the most recent window whose step $k$
has already been observed, so the procedure uses nothing an operator would not have (the
immediate-feedback form, which sees the whole horizon one hour later, is kept in the result
files as a diagnostic: it reaches the same coverage, within
\nv{prob.grid.lookahead_gap_pp} points on average, with intervals that differ by
\nv{prob.grid.lookahead_width_pct} in width and by \nv{prob.grid.lookahead_winkler_pct} in
Winkler score, which is the optimism the delayed form removes). The scores are the mean pinball loss over the seven levels; the quantile-weighted
CRPS approximation (qCRPS), the trapezoidal integral of twice the pinball loss over the emitted
levels, which does not score the tails beyond the outer levels and is not the CRPS of a
continuous predictive distribution; and, for the 90\,\% interval, coverage, mean width and the
Winkler score. Differences in sharpness therefore reflect the forecaster rather than the
calibrator.

\subsection{Reproducibility}
Every table, figure and number in this paper is regenerated from the released result files by
three scripts and one build command; the prose contains no hand-typed result, which is
enforced by a check in the build. The released decision log records every protocol decision
taken during development and every change to one, with the validation evidence that motivated
it. Every neural and zero-shot run used a single 8\,GB Quadro RTX 4000 (driver 596.51) with
Python 3.11.9, PyTorch 2.4.1 (CUDA 12.1, cuDNN 9.1), NeuralForecast 3.1.8, NumPy 1.26.4 and
scikit-learn 1.8.0; the classical references ran on CPU. Seeds fix the initialisation, the
window order and the weather-noise stream; cuDNN kernels are not forced deterministic, so a
re-run reproduces each cell to within the seed spread printed beside it rather than bit for bit.
Every result record carries its full configuration.
